\documentclass[11pt]{article}
\usepackage[margin=1in]{geometry}
\usepackage{amsmath,amssymb}
\usepackage[hidelinks]{hyperref}

\begin{document}

\begin{center}
{\Large Property \(\mathcal A\) does not force frequent hypercyclicity}\\[4pt]
{\normalsize Literature-implied status packet for arXiv:1008.4017}
\end{center}

\section*{Status}

This is a \texttt{literature\_implied\_answer} packet. It records that
Menet's later chaotic non-frequently-hypercyclic operator gives a negative
answer to a stronger question posed by Costakis--Parissis in
\emph{Szemeredi's theorem, frequent hypercyclicity and multiple recurrence},
arXiv:1008.4017 \cite{CP}.

\section*{Source question}

In Section 2, Costakis--Parissis define property \(\mathcal A\): a sequence
\((T_n)_{n\in\mathbb N}\) has property \(\mathcal A\) if for every open
set \(U\subset X\) there is \(x\in X\) such that
\[
  \overline{\mathrm{dens}}\{n\in\mathbb N:T_nx\in U\}>0.
\]
For a single operator \(T\), this means that \((T^n)_{n\in\mathbb N}\) has
property \(\mathcal A\).

In the final ``Further Questions'' section they write:
\begin{quote}
It is easy to see that every chaotic operator \(T\) has property
\(\mathcal A\). A well known question asks whether every chaotic operator is
frequently hypercyclic; see for example \cite{BM}. An even stronger question
is thus whether every hypercyclic operator that has property \(\mathcal A\)
is frequently hypercyclic.
\end{quote}

\section*{Supporting answer}

Quentin Menet's later paper \emph{Linear chaos and frequent hypercyclicity},
arXiv:1410.7173 \cite{Menet}, answers the embedded chaotic-operator problem
negatively. Theorem 1.2 states that there exists a chaotic operator \(T\) on
\(\ell^1\) which is neither \(\mathcal U\)-frequently hypercyclic nor
distributionally chaotic; in particular, \(T\) is chaotic and not frequently
hypercyclic.

Combining these two facts gives the advertised counterexample to the stronger
source question. Menet's operator is chaotic, so it is hypercyclic. By the
Costakis--Parissis observation, it has property \(\mathcal A\). But Theorem
1.2 of Menet shows that it is not frequently hypercyclic. Therefore
\[
  \text{hypercyclic} + \mathcal A \;\not\Rightarrow\;
  \text{frequently hypercyclic}.
\]

\section*{Scope limitations}

This packet is not a new construction from the run; it is a short
later-literature implication. It settles only the frequent-hypercyclicity
question above. The subsequent Costakis--Parissis question asking whether
there exists an operator with property \(\mathcal A\) which is not
\(\mathcal U\)-frequently recurrent is not settled here, because Menet's
Theorem 1.2 concerns \(\mathcal U\)-frequent hypercyclicity rather than
\(\mathcal U\)-frequent recurrence.

\begin{thebibliography}{9}
\bibitem{CP}
G. Costakis and I. Parissis,
\emph{Szemeredi's theorem, frequent hypercyclicity and multiple recurrence},
arXiv:1008.4017.

\bibitem{BM}
F. Bayart and E. Matheron,
\emph{Dynamics of linear operators},
Cambridge University Press, 2009.

\bibitem{Menet}
Q. Menet,
\emph{Linear chaos and frequent hypercyclicity},
arXiv:1410.7173.
\end{thebibliography}

\end{document}
